\chapter{Circuit Modeling: Paradigms and Components}
\label{ch:foundations}
Part~II built a generic linear-systems toolkit. Part~III
now grounds it in real circuits. Before any specific circuit, this chapter sets up
the modeling paradigm the rest of the book uses---what a ``circuit'' even
is---then introduces the ideal components (\(R\), \(L\), \(C\), and sources), the
two conservation laws that connect them, the state variables that carry a
circuit's memory, and the unit bookkeeping that catches errors. Everything here is
the vocabulary the later chapters use.

\section{The Lumped-Element Modeling Paradigm}
A real circuit is a tangle of fields spread through space. To analyze it we make
the \textbf{lumped-element idealization}: we pretend all the interesting physics
lives inside a few discrete two-terminal components joined by perfect wires. Each
component is described by a single relationship between the voltage across it and
the current through it; the wires carry current with no voltage drop and store no
energy. This is reasonable whenever the circuit is small compared with the wavelength of
the signals in it, so that a change at one end is felt essentially instantly at the
other. When that assumption fails---long feed lines, antennas at radio
frequencies (RF)---we switch to
the distributed model of \cref{ch:lines}. Until then, ``the circuit'' means an
idealized network of lumped elements.

\begin{controlsbox}
Lumping is the electrical version of the modeling step every engineer makes:
replace a spatially distributed reality (a heat-conducting bar, a flexing beam,
a mixing tank) with a finite set of state variables governed by ordinary---not
partial---differential equations. The lumped circuit \emph{is} the ODE model.
\end{controlsbox}

\section{The Ideal Components}
Four idealized elements carry essentially all of the exam's analog circuit
content. Each is defined by its \textbf{element law} (or constitutive relation)
relating its voltage and current:
\begin{description}
\item[Resistor \(R\)] \(v_R=Ri_R\) (Ohm's law). Dissipates energy as heat;
carries no memory. Unit: ohm, \(\si{\ohm}\).
\item[Capacitor \(C\)] \(i_C=C\,\dd v_C/\dd t\). Stores energy in an electric
field; its voltage cannot jump under finite current. Unit: farad, \(\si{\farad}\).
\item[Inductor \(L\)] \(v_L=L\,\dd i_L/\dd t\). Stores energy in a magnetic
field; its current cannot jump under finite voltage. Unit: henry, \(\si{\henry}\).
\item[Sources] an ideal voltage source fixes its terminal voltage regardless of
current; an ideal current source fixes its current regardless of voltage. Sources
are how energy and signals enter the network.
\end{description}
The two storage laws are the only places a derivative enters, which is why
capacitors and inductors---not resistors---are what give a circuit dynamics. The
stored energies are
\[
E_C=\tfrac{1}{2}Cv_C^2,\qquad
E_L=\tfrac{1}{2}Li_L^2.
\]

\begin{physicalbox}
An ideal component is a modeling fiction: a real resistor has a little
inductance, a real capacitor a little leakage. The idealization keeps the model
tractable, and the corrections are added back only when they matter (for example,
a capacitor's series inductance at very high frequencies, VHF).
\end{physicalbox}

\subsection{Self-Resonance: When a Part Becomes Its Opposite}
\label{sec:selfresonance}
That last point deserves a number rather than a caveat, because it decides which
component you can use at which frequency.

Every real component carries a parasitic reactance of the \emph{other} kind. A
capacitor's leads and plates have inductance \(L_s\) in series with \(C\); a coil's
adjacent turns have capacitance \(C_p\) in parallel with \(L\). Each pairing is
therefore an \(L\)--\(C\) resonator, and it resonates at the frequency any \(LC\)
pair does:
\[
f_{\mathrm{SRF}}=\frac{1}{2\pi\sqrt{L_sC}}
\quad\text{(capacitor)},
\qquad
f_{\mathrm{SRF}}=\frac{1}{2\pi\sqrt{LC_p}}
\quad\text{(inductor)} .
\]
This is the \textbf{self-resonant frequency}, and the consequence is worth stating
sharply. Below it the part behaves as marked. Above it the \emph{parasitic}
dominates, so
\begin{itemize}
\item a capacitor above its \(f_{\mathrm{SRF}}\) is an \emph{inductor}---its
impedance rises with frequency instead of falling, which is why a bypass capacitor
stops bypassing;
\item an inductor above its \(f_{\mathrm{SRF}}\) is a \emph{capacitor}---its
impedance falls with frequency, which is why an RF choke stops choking.
\end{itemize}
A capacitor is a series resonator, so it is a near-short at \(f_{\mathrm{SRF}}\); an
inductor is a parallel resonator, so it is a near-open there. Those are exactly the
two behaviors of \cref{ch:rlc,ch:rlcpar}, which is why this subsection needs no new
theory---only the observation that the resonator was there all along.

\begin{workedbox}[: where a bypass capacitor quits]
A \(\SI{10}{\nano\farad}\) ceramic capacitor with \SI{10}{\milli\meter} of total lead
and plate length carries roughly \(L_s\approx\SI{10}{\nano\henry}\)
(about \SI{1}{\nano\henry} per millimetre). Then
\[
f_{\mathrm{SRF}}=\frac{1}{2\pi\sqrt{(10\times10^{-9})(10\times10^{-9})}}
=\frac{1}{2\pi(10^{-8})}\approx\SI{15.9}{\mega\hertz}.
\]
So this part is a good bypass through the high-frequency (HF) bands and useless at
VHF---at
\SI{144}{\mega\hertz} its \SI{10}{\nano\henry} presents \(\omega L_s\approx\SI{9}{\ohm}\)
of \emph{inductive} reactance while the capacitance it was bought for would have
offered only \(\SI{0.11}{\ohm}\). The standard fix follows from the formula:
shorten the leads (lower \(L_s\)), or parallel a large capacitor with a small one
whose \(f_{\mathrm{SRF}}\) is higher.
\end{workedbox}

\section{KCL as an Accumulation Balance}
Kirchhoff's current law (KCL) is a charge balance. At a node,
\[
\text{accumulation}=\text{inflow}-\text{outflow}.
\]
For an ordinary wire node with negligible stored charge, accumulation is zero,
so the algebraic sum of currents is zero:
\[
\sum_k i_k=0.
\]

\begin{controlsbox}
This is the same structure as a material balance in process systems. Capacitors
make the accumulation term visible because \(i_C=C\,\dd v_C/\dd t\).
\end{controlsbox}

\section{KVL as Energy Consistency}
Kirchhoff's voltage law (KVL) says that the signed voltage rises and drops around a
closed loop sum to zero:
\[
\sum_k v_k=0.
\]
Voltage is energy per unit charge. KVL is therefore an energy-consistency rule:
after a charge goes around a closed loop and returns to the same point, its net
change in potential energy is zero.

\section{State Variables}
The natural state variables are the quantities attached to stored energy:
capacitor voltages and inductor currents. A circuit with one independent energy
storage element is usually first order. A circuit with two independent storage
elements is usually second order. These are exactly the variables that cannot jump,
so they are what the system ``remembers''---the electrical instance of the state
vector of \cref{ch:linsys}.

\section{Electrical Units and Dimensional Checks}
Units are not decoration; they are a fast way to catch a wrong formula. The base
relationships all follow from the definitions above:
\[
\SI{1}{\coulomb}=\SI{1}{\ampere\second},
\qquad
\SI{1}{\volt}=\SI{1}{\joule\per\coulomb},
\qquad
\SI{1}{\ohm}=\SI{1}{\volt\per\ampere}.
\]
Note the direction of these definitions: the \textbf{ampere is an SI base unit},
defined by fixing the elementary charge, so current is the primitive electrical
quantity and the coulomb is \emph{derived} from it (\(\SI{1}{\coulomb}
=\SI{1}{\ampere\second}\)), not the other way around. Everything else here---volt,
ohm, farad, henry, watt---is then built from the ampere together with the mechanical
base units.
\[
\SI{1}{\farad}=\SI{1}{\coulomb\per\volt},
\qquad
\SI{1}{\henry}=\SI{1}{\volt\second\per\ampere},
\qquad
\SI{1}{\watt}=\SI{1}{\joule\per\second}=\SI{1}{\volt\ampere}.
\]
Hertz counts cycles per second; radians per second is the angular rate used in
calculus:
\[
\omega=2\pi f.
\]

\subsection{Radio Prefixes}
(\Cref{app:units} repeats this table, and the unit definitions above, as a one-page
reference.)
\begin{tabularx}{\textwidth}{@{}L{0.2\textwidth}L{0.2\textwidth}Y@{}}
\toprule
Prefix & Factor & Example \\
\midrule
pico & \(10^{-12}\) & \(\SI{220}{\pico\farad}\) coupling capacitor. \\
nano & \(10^{-9}\) & \(\SI{10}{\nano\farad}\) bypass capacitor. \\
micro & \(10^{-6}\) & \(\SI{10}{\micro\henry}\) RF inductor. \\
milli & \(10^{-3}\) & \(\SI{25}{\milli\ampere}\) current. \\
kilo & \(10^{3}\) & \(\SI{4.7}{\kilo\ohm}\) resistor. \\
mega & \(10^{6}\) & \(\SI{7.1}{\mega\hertz}\) frequency. \\
giga & \(10^{9}\) & \(\SI{1.2}{\giga\hertz}\) microwave frequency. \\
\bottomrule
\end{tabularx}

\subsection{Useful Dimensional Tests}
\begin{unitsbox}
\[
RC=(\si{\ohm})(\si{\farad})
=\left(\frac{\si{\volt}}{\si{\ampere}}\right)
\left(\frac{\si{\ampere\second}}{\si{\volt}}\right)
=\si{\second}.
\]
\[
\frac{L}{R}
=\frac{\si{\volt\second\per\ampere}}{\si{\volt\per\ampere}}
=\si{\second}.
\]
\[
LC=(\si{\henry})(\si{\farad})=\si{\second\squared},
\qquad
\frac{1}{\sqrt{LC}}=\si{\per\second}.
\]
\[
\omega L=(\si{\per\second})(\si{\volt\second\per\ampere})
=\si{\ohm},
\qquad
\frac{1}{\omega C}=\si{\ohm}.
\]
\end{unitsbox}

\subsection{Worked Unit Check}
Suppose \(R=\SI{10}{\kilo\ohm}\) and \(C=\SI{220}{\pico\farad}\). Then
\[
\tau=RC=(10\times10^3)(220\times10^{-12})=\SI{2.2e-6}{\second}.
\]
The corner (cutoff) angular frequency is the reciprocal of the time constant,
\[
\omega_c=\frac{1}{RC}=\frac{1}{\tau}=\SI{4.55e5}{\radian\per\second},
\qquad
f_c=\frac{\omega_c}{2\pi}\approx \SI{72.3}{\kilo\hertz}.
\]

\begin{mistakebox}
Do not use \(f\) where an equation requires \(\omega\). Reactance formulas use
\(\omega\): \(X_L=\omega L\) and \(X_C=1/(\omega C)\). If frequency is given in
hertz, convert with \(\omega=2\pi f\).
\end{mistakebox}

\begin{exambox}
Nearly every component question on the exam is really a question about which
idealization has failed. An ideal capacitor blocks DC and passes RF; a real one has
lead inductance and eventually resonates. An ideal inductor passes DC and blocks
RF; a real one has winding capacitance and does the same. When a question describes
a part behaving as its \emph{opposite} type, it is telling you the frequency is
above that part's self-resonance (\cref{sec:selfresonance}). When it describes a
wire or a ground lead mattering, the lumped assumption has failed and
\cref{ch:lines} applies.
\end{exambox}

\section{Summary: Key Equations and Definitions}
\[
\sum i_k=0,\qquad \sum v_k=0,\qquad
p(t)=v(t)i(t).
\]

\medskip
\begin{description}
\item[Lumped element] an idealized two-terminal component whose behavior is a
single \(v\)--\(i\) law; valid when the circuit is small versus a wavelength.
\item[KCL] Conservation of charge at a node.
\item[KVL] Energy consistency around a loop.
\item[State] A variable needed, with the input, to predict future behavior
(capacitor voltages, inductor currents).
\item[Passive sign convention] Positive absorbed power when current enters the
positive-voltage terminal.
\end{description}
